\section*{NeurIPS Paper Checklist}

\begin{enumerate}

\item {\bf Claims}
    \item[] Question: Do the main claims made in the abstract and introduction accurately reflect the paper's contributions and scope?
    \item[] Answer: \answerYes{}
    \item[] Justification: The abstract and introduction state the discrete-action self-play setting, the adversarial action-removal threat model, the empirical scope, and the CACw/CACv mechanism. The Limitations section further bounds the claims to discrete action spaces and empirical scaling trends.

\item {\bf Limitations}
    \item[] Question: Does the paper discuss the limitations of the work performed by the authors?
    \item[] Answer: \answerYes{}
    \item[] Justification: The paper includes a Limitations section discussing discrete-action scope, local theory, independence assumptions, and remaining full-scale poker/continuous-action gaps.

\item {\bf Theory assumptions and proofs}
    \item[] Question: For each theoretical result, does the paper provide the full set of assumptions and a complete (and correct) proof?
    \item[] Answer: \answerYes{}
    \item[] Justification: Propositions state their assumptions, and the paper gives proof sketches plus an explicit theoretical-scope paragraph. The results are intentionally local/sufficient rather than a global optimal-adversary characterization.

\item {\bf Experimental result reproducibility}
    \item[] Question: Does the paper fully disclose all the information needed to reproduce the main experimental results of the paper to the extent that it affects the main claims and/or conclusions of the paper?
    \item[] Answer: \answerYes{}
    \item[] Justification: The Reproducibility Statement identifies the experiment scripts, fixed seeds, and figure-generation script. Hyperparameters and normalization bounds appear in the appendix.

\item {\bf Open access to data and code}
    \item[] Question: Does the paper provide open access to the data and code, with sufficient instructions to faithfully reproduce the main experimental results, as described in supplemental material?
    \item[] Answer: \answerYes{}
    \item[] Justification: The repository includes all environments, agents, adversaries, and scripts. No external datasets are used. For double-blind submission, the repository should be anonymized before linking.

\item {\bf Experimental setting/details}
    \item[] Question: Does the paper specify all the training and test details necessary to understand the results?
    \item[] Answer: \answerYes{}
    \item[] Justification: The paper describes self-play, masking, adversary training, action-removal granularity, DQN masking mechanics, and main hyperparameters. Additional details are in Appendix~\ref{app:hyperparams}.

\item {\bf Experiment statistical significance}
    \item[] Question: Does the paper report error bars suitably and correctly defined or other appropriate information about the statistical significance of the experiments?
    \item[] Answer: \answerYes{}
    \item[] Justification: Main experiments report 95\% confidence intervals over seeds where applicable. Some auxiliary diagnostic tables report means and are described as diagnostics.

\item {\bf Experiments compute resources}
    \item[] Question: For each experiment, does the paper provide sufficient information on the computer resources needed to reproduce the experiments?
    \item[] Answer: \answerYes{}
    \item[] Justification: Appendix tables report representative sample counts and wall-clock times. Experiments run on CPU with PyTorch used for neural agents and adversaries.

\item {\bf Code of ethics}
    \item[] Question: Does the research conducted in the paper conform, in every respect, with the NeurIPS Code of Ethics?
    \item[] Answer: \answerYes{}
    \item[] Justification: The work uses synthetic environments and does not involve human subjects, private data, or deployed systems. It is framed as robustness/security research with defensive implications.

\item {\bf Broader impacts}
    \item[] Question: Does the paper discuss both potential positive societal impacts and negative societal impacts of the work performed?
    \item[] Answer: \answerYes{}
    \item[] Justification: The paper discusses capability-removal risks in APIs, actuators, and platforms, and motivates defenses that preserve high-CACv decision capacity. The work could be misused to identify targeted capability restrictions, but it is presented to motivate robustness testing and defenses.

\item {\bf Safeguards}
    \item[] Question: Does the paper describe safeguards that have been put in place for responsible release of data or models that have a high risk for misuse?
    \item[] Answer: \answerNA{}
    \item[] Justification: No pretrained models, scraped datasets, or high-risk generative artifacts are released. The code contains small synthetic RL environments and experiment scripts.

\item {\bf Licenses for existing assets}
    \item[] Question: Are the creators or original owners of assets used in the paper properly credited and are the license and terms of use explicitly mentioned and properly respected?
    \item[] Answer: \answerYes{}
    \item[] Justification: The code uses standard open-source Python packages and cites related papers. No external datasets are used. A repository license should be added before broad public release.

\item {\bf New assets}
    \item[] Question: Are new assets introduced in the paper well documented and is the documentation provided alongside the assets?
    \item[] Answer: \answerYes{}
    \item[] Justification: New assets are code, synthetic environments, and paper sources. The README documents layout, dependencies, run commands, scope, and reproducibility.

\item {\bf Crowdsourcing and research with human subjects}
    \item[] Question: For crowdsourcing experiments and research with human subjects, does the paper include the full text of instructions and compensation details?
    \item[] Answer: \answerNA{}
    \item[] Justification: No crowdsourcing or human-subject experiments are used.

\item {\bf Institutional review board (IRB) approvals or equivalent for research with human subjects}
    \item[] Question: Does the paper describe potential participant risks and IRB approvals where applicable?
    \item[] Answer: \answerNA{}
    \item[] Justification: The work does not involve human subjects.

\item {\bf Declaration of LLM usage}
    \item[] Question: Does the paper describe the usage of LLMs if it is an important, original, or non-standard component of the core methods in this research?
    \item[] Answer: \answerNA{}
    \item[] Justification: The core methodology, experiments, and analysis do not use LLMs.

\end{enumerate}
